\scanpage{44}\mspage{37}
% archive sheet number: 43
\begingroup\makeatletter\tagsleft@true\makeatother
Les inclusions de (iii).

L'inclusion
\[
 C\subset\operatorname{Cof}_{W}
\]
implique que toute $f\in C(I)$ est dans $\operatorname{Cof}_{W}$ argument par argument. Cela \unsure{indique|implique|veut que|veut dire} que $f$ est dans $\operatorname{Cof}_{W(I)}$, comme on voit aussitôt sur la définition $\leftarrow$ je l'explicite en le

\uline{Lemme 1} : Soit $f$ ou $X_{\unsure{\cdot}}\unsure{-|\to}Y_{\cdot}$ une flèche de $M(I)$. Pour que $f$ soit dans $\operatorname{Cof}_{W(I)}$ (resp. dans $\operatorname{Fib}_{W(I)}$), il \uline{suffit} que pour tt $s$ dans $\operatorname{Ob} I$, $\alpha_{s}^{*}(f)$ soit dans $\operatorname{Cof}_{W}$ (resp. dans $\operatorname{Fib}_{W}$). Idem (pour mémoire) pour \unsure{la} condition pour que \deleted{\unsure{\(f\)}} $f\in W(I)$ -- là il \uline{f et s} que $\alpha_{s}^{*}(f)\in W\quad\forall s\in\operatorname{Ob}I$.

Nous allons appliquer ce $=$ lemme pour \deleted{\unsure{ét}} établir les inclusions
\[
 C(I)_{*}\subset W(I),\quad (C(I)\cap W(I))_{*}\subset\operatorname{Fib}_{\deleted{\unsure{I}}W}(I)
\]
Donc il suffit de voir que
\[
\begin{aligned}
 &\alpha_{s}^{*}(C(I)_{*})\subset C_{*}\quad(\deleted{\unsure{\alpha_s}}\subset W)\\
 &\deleted{\unsure{\alpha_{s}^{*}((C(I)\cap W)_{*}))\subset}}\\
 &\alpha_{s}^{*}((\underbrace{C(I)\cap W(I)}_{(C\cap W)(I)})_{*})\subset(C\cap W)_{*}\quad(\subset\operatorname{Fib}_{W}),
\end{aligned}
\]
et plus généralement
\[
 \alpha_{s}^{*}(\unsure{\Phi|\mathscr P}(I)_{*})\subset\Phi_{*}\tag*{(7)}
\]
pour tte \inserted{(!)} partie $\Phi$ de $\operatorname{Fl}(M)$, définissant donc $\Phi(I)\subset\operatorname{Fl}(M(I))$ terme à terme.
\endgroup
